\chapter{Long-Haul Wireless}\label{ch:nw:long-haul-wireless}

When a storm crosses the line of towers between Chicago and New Jersey, the \omterm{def:nw:long-haul-wireless:mw}{microwave links} under it fade, and for as long as the rain
lasts the fastest route between the futures and the stocks is fibre again, three milliseconds slower. Shkilko and Sokolov used exactly
those afternoons as an experiment: weather episodes that ``temporarily remove the speed advantages of the fastest traders by disrupting
their microwave networks'' were associated with lower adverse selection and lower trading costs. The rain gave the slower traders a few
hours without the race, and the market's liquidity improved while it lasted.

Chapter~10 measured the routes against their floors and chapter~13 took the glass apart. This chapter does the same for the air: the geometry
of a line-of-sight hop, the power it has to spare, the rain that takes it away, the millimetre-wave and laser links inside a metropolitan
area, the shortwave radio that crosses oceans, and what these links can carry. Its model runs a labelled storm across a chain of towers
between Aurora and Carteret and counts the minutes on fibre.

\section{Microwave links: line of sight, Fresnel zones and towers}

\begin{definition}[Microwave link, line-of-sight path]\label{def:nw:long-haul-wireless:mw}
A \emph{microwave link}\index{microwave link} carries data by radio at frequencies of a few to a few tens of gigahertz between two directional
antennas, typically on towers or rooftops. It needs a \emph{line-of-sight path}\index{line-of-sight path}: a straight path between the antennas
clear of the ground and of obstacles, since such waves do not bend around them.
\end{definition}

\begin{definition}[Fresnel zone, earth bulge]\label{def:nw:long-haul-wireless:fresnel}
The first \emph{Fresnel zone}\index{Fresnel zone} of a radio path is the ellipsoid around the straight line between the antennas in which
paths are at most half a wavelength longer than the direct one; obstacles inside it weaken the received signal, so a link is designed to keep
it clear. The \emph{earth bulge}\index{earth bulge} is the height of the Earth's surface above the straight line joining two points at ground
level, at a given point between them; atmospheric refraction bends radio waves slightly downward, which is modelled by an earth of radius
$K$ times the real one.
\end{definition}

\begin{proposition}[Hop geometry]\label{prop:nw:long-haul-wireless:geometry}
At distances $d_1$ and $d_2$ kilometres from the two ends of a hop of length $d = d_1 + d_2$, the \omterm{def:nw:long-haul-wireless:fresnel}{earth bulge} is
$h = d_1 d_2 / (12.74\,K)$ metres and the radius of the first \omterm{def:nw:long-haul-wireless:fresnel}{Fresnel zone} at $f$ gigahertz is $r = 17.32\sqrt{d_1 d_2 / (f\,d)}$ metres. At
mid-hop, $h = d^2/(50.96\,K)$ and $r = 8.66\sqrt{d/f}$: the bulge grows with the square of the hop, the zone with its square root.
\end{proposition}

\begin{proof}
A circle of radius $K R_\oplus$ rises above its chord, at distances $d_1$ and $d_2$ from its ends, by $d_1 d_2/(2 K R_\oplus)$ to first order;
with $R_\oplus = 6\,371$~km and the result in metres, $2 R_\oplus/1\,000 = 12.74$. The first Fresnel radius is $\sqrt{\lambda d_1 d_2/d}$ with
$\lambda = c_0/f$; with $d$ in kilometres and $f$ in gigahertz the constant is $\sqrt{0.2998 \times 10^3} = 17.32$. Setting $d_1 = d_2 = d/2$
gives the mid-hop forms.
\end{proof}

Laughlin, Aguirre and Grundfest give the same mid-hop rules (8.7 and 1/50) and draw the consequence for the Chicago--New York networks: with
towers of about 100 metres, frequencies near \qty{6}{\giga\hertz}, $K$ below 4/3 and obstacles of 10 metres, a hop over flat terrain cannot
be much longer than about 70 kilometres, so a 1\,200-kilometre route needs twenty hops or more. \Cref{tab:nw:long-haul-wireless:geometry}
reproduces the arithmetic.

\begin{table}[ht]
\centering\small
\begin{tabular}{rrrrr}
\toprule
Hop & Frequency & \omterm{def:nw:long-haul-wireless:fresnel}{Earth bulge} & Fresnel radius & Tower height \\
(km) & (GHz) & at mid-hop (m) & at mid-hop (m) & (m) \\
\midrule
30 & 6 & 13.2 & 19.4 & 42.6 \\
30 & 11 & 13.2 & 14.3 & 37.5 \\
50 & 6 & 36.8 & 25.0 & 71.8 \\
50 & 11 & 36.8 & 18.5 & 65.3 \\
70 & 6 & 72.1 & 29.6 & 111.7 \\
70 & 11 & 72.1 & 21.8 & 104.0 \\
\bottomrule
\end{tabular}
\caption{Hop geometry from \cref{prop:nw:long-haul-wireless:geometry} with $K = 4/3$: the height two equal towers need so that a 10-metre obstacle
at mid-hop stays clear of the whole first \omterm{def:nw:long-haul-wireless:fresnel}{Fresnel zone}. Beyond 70 kilometres the towers pass 100 metres. Data:
\texttt{nw\_wireless.geometry\_rows()}.}\label{tab:nw:long-haul-wireless:geometry}
\end{table}

\begin{omfigure}
\begin{tikzpicture}[font=\footnotesize,>=Stealth]
\draw[thick,gray] (0,0) arc[start angle=110,end angle=70,radius=16.4];
\draw[thick] (0,0) -- (0,2.3); \draw[thick] (11.2,0) -- (11.2,2.3);
\node[below] at (0,0) {tower A}; \node[below] at (11.2,0) {tower B};
\draw[omDef,thick] (0,2.3) -- (11.2,2.3);
\draw[omThm,dashed] (5.6,2.3) ellipse (5.6 and 0.55);
\draw[gray,densely dotted] (0,0) -- (11.2,0);
\draw[<->,omMeth] (5.6,0) -- node[right]{earth bulge} (5.6,0.99);
\draw[<->,omThm] (6.8,2.3) -- node[right]{Fresnel radius} (6.8,1.78);
\node[omDef] at (3.0,3.1) {line of sight};
\fill[gray!50] (4.9,0.99) rectangle (5.3,1.3); \node[font=\scriptsize] at (4.4,1.2) {obstacle};
\end{tikzpicture}

\omcaption{A microwave hop, schematically: the towers must lift the line of sight above the \omterm{def:nw:long-haul-wireless:fresnel}{earth bulge} at mid-hop, plus the obstacles, plus the
first \omterm{def:nw:long-haul-wireless:fresnel}{Fresnel zone} around the line. Vertical scale exaggerated.}\label{fig:nw:long-haul-wireless:hop}
\end{omfigure}

\begin{definition}[Link budget, fade margin]\label{def:nw:long-haul-wireless:budget}
A link's \emph{link budget}\index{link budget} adds the transmitter's power and the antennas' gains and subtracts the path's losses to give the
power received. Its \emph{fade margin}\index{fade margin} is the received power minus the least power at which the receiver still decodes the
signal at the required error rate: the loss the link can absorb before it fails.
\end{definition}

The budget follows from physics once the equipment is fixed. Free space loses $92.45 + 20\log_{10} f + 20\log_{10} d$ decibels ($f$ in
gigahertz, $d$ in kilometres); a dish of diameter $D$ gains $10\log_{10}[\eta(\pi D f/c_0)^2]$. With the model's equipment (30~dBm
transmitted, 1.8-metre dishes of 60\% efficiency, \qty{3}{\decibel} of other losses, a receiver threshold of $-70$~dBm, all assumptions),
a 60-kilometre hop keeps between 31 and \qty{54}{\decibel} of margin (\cref{tab:nw:long-haul-wireless:budget}). Higher frequencies gain more in
the dishes than they lose in free space; the rain takes it back.

\section{Weather, rain fade and availability}

\begin{definition}[Rain fade, link availability]\label{def:nw:long-haul-wireless:rain}
\emph{Rain fade}\index{rain fade} is the attenuation of a radio signal by raindrops along its path, which grows with the rain rate and steeply
with frequency. A link's \emph{link availability}\index{link availability} is the fraction of time in which it carries traffic at its required
error rate, usually stated per year.
\end{definition}

Recommendation ITU-R P.838-3 gives rain's specific attenuation as a power law of the rain rate, $\gamma = k R^{\alpha}$ decibels per kilometre
for $R$ millimetres an hour, with coefficients tabulated by frequency and polarisation. At \qty{11}{\giga\hertz} (horizontal polarisation)
$k = 0.01772$ and $\alpha = 1.2140$: a storm of \qty{50}{\milli\metre\per\hour} costs \qty{2.05}{\decibel} per kilometre. At \qty{6}{\giga\hertz},
$k = 0.0007056$ and $\alpha = 1.59$; at \qty{80}{\giga\hertz}, $k = 1.1704$ and $\alpha = 0.7115$.
\Cref{fig:nw:long-haul-wireless:gamma} shows the spread: in a light rain of \qty{2.5}{\milli\metre\per\hour}, a factor of about 740 between 6 and \qty{80}{\giga\hertz}.

\begin{omfigure}
\begin{tikzpicture}
\begin{semilogyaxis}[width=11cm,height=6.2cm,xmin=0,xmax=150,ymin=0.0005,ymax=60,xlabel={rain rate (mm/h)},
  ylabel={specific attenuation (dB/km)},tick label style={font=\footnotesize},label style={font=\small},grid=major,
  grid style={gray!25},legend pos=south east,legend style={font=\footnotesize},table/col sep=comma]
\addplot[omDef,thick,mark=*,mark size=1.3pt] table[x=rain,y=g6]{figdata/networks/14-long-haul-wireless/gamma.csv};
\addplot[omThm,thick,mark=square*,mark size=1.3pt] table[x=rain,y=g11]{figdata/networks/14-long-haul-wireless/gamma.csv};
\addplot[omMeth,thick,mark=triangle*,mark size=1.6pt] table[x=rain,y=g18]{figdata/networks/14-long-haul-wireless/gamma.csv};
\addplot[gray,thick,mark=diamond*,mark size=1.6pt] table[x=rain,y=g80]{figdata/networks/14-long-haul-wireless/gamma.csv};
\legend{6 GHz,11 GHz,18 GHz,80 GHz}
\end{semilogyaxis}
\end{tikzpicture}

\omcaption{Specific attenuation of rain against rain rate at four frequencies, horizontal polarisation, from the ITU-R P.838-3 coefficients. Data:
\texttt{fig\_wireless.py}; coefficients in \texttt{code/firm/radiolink/data/p838\_3.csv}.}\label{fig:nw:long-haul-wireless:gamma}
\end{omfigure}

\begin{proposition}[The rain that breaks a hop]\label{prop:nw:long-haul-wireless:critical}
A hop with \omterm{def:nw:long-haul-wireless:budget}{fade margin} $M$ decibels fails when rain of rate $R$ covers a length $L$ of it with $k R^{\alpha} L > M$, that is above the critical
rate
\[
  R^{*} = \left(\frac{M}{k L}\right)^{1/\alpha}.
\]
\end{proposition}

\begin{proof}
The rain's attenuation over the wet length is $\gamma L = k R^{\alpha} L$; the hop fails when it exceeds the margin. Solving for $R$ gives
$R^{*}$.
\end{proof}

A 60-kilometre hop at \qty{11}{\giga\hertz} keeps \qty{36.4}{\decibel} of margin in the model: rain of \qty{18.4}{\milli\metre\per\hour} along
its whole length would break it, but a storm cell ten kilometres across needs \qty{80}{\milli\metre\per\hour}. At \qty{6}{\giga\hertz} the same
hop survives \qty{64}{\milli\metre\per\hour} everywhere; at \qty{80}{\giga\hertz} it fails in a drizzle, which is why long hops use the lower
bands and short ones can use the higher.

\begin{table}[ht]
\centering\small
\begin{tabular}{rrrrrrr}
\toprule
Frequency & Dish gain & Free-space loss & Received & \omterm{def:nw:long-haul-wireless:budget}{Fade margin} & $R^*$, whole hop & $R^*$, 10-km cell \\
(GHz) & (dBi) & (dB) & (dBm) & (dB) & (mm/h) & (mm/h) \\
\midrule
6 & 38.9 & 143.6 & $-38.9$ & 31.1 & 63.5 & 196.0 \\
11 & 44.1 & 148.8 & $-33.6$ & 36.4 & 18.4 & 80.4 \\
18 & 48.4 & 153.1 & $-29.3$ & 40.7 & 8.1 & 42.3 \\
80 & 61.4 & 166.1 & $-16.4$ & 53.6 & 0.7 & 8.5 \\
\bottomrule
\end{tabular}
\caption{\omterm{def:nw:long-haul-wireless:budget}{Link budget} of a 60-kilometre hop with the model's equipment (30~dBm, two 1.8-metre dishes at 60\% efficiency, \qty{3}{\decibel}
of other losses, $-70$~dBm threshold), and the critical rain rate of \cref{prop:nw:long-haul-wireless:critical} for rain along the whole hop
and for a cell ten kilometres wide. Data: \texttt{nw\_wireless.budget\_rows()}.}\label{tab:nw:long-haul-wireless:budget}
\end{table}

\omcode{code/firm/radiolink/firm_radiolink.py}{83}{97}{Rain on a hop: the ITU-R power law over the wet length, and the rate at which it consumes
the fade margin.}\label{lst:nw:long-haul-wireless:rain}

\begin{omfigure}
\begin{tikzpicture}
\begin{axis}[width=11cm,height=6cm,xmin=5,xmax=80,ymin=0,ymax=440,xlabel={hop length (km)},
  ylabel={critical rain rate (mm/h)},tick label style={font=\footnotesize},label style={font=\small},grid=major,
  grid style={gray!25},legend columns=3,legend style={font=\footnotesize,at={(0.5,-0.25)},anchor=north,draw=none},table/col sep=comma]
\addplot[omDef,thick,no markers] table[x=d_km,y=r6]{figdata/networks/14-long-haul-wireless/critical.csv};
\addplot[omThm,thick,no markers] table[x=d_km,y=r11]{figdata/networks/14-long-haul-wireless/critical.csv};
\addplot[omMeth,thick,no markers] table[x=d_km,y=r18]{figdata/networks/14-long-haul-wireless/critical.csv};
\legend{6 GHz,11 GHz,18 GHz}
\end{axis}
\end{tikzpicture}

\omcaption{The rain rate at which a hop loses its \omterm{def:nw:long-haul-wireless:budget}{fade margin}, against hop length, for a rain cell ten kilometres across the hop (the whole hop,
below ten kilometres), with the model's equipment. Longer hops lose margin to free space; higher frequencies lose it to rain. Data: \texttt{fig\_wireless.py},
\texttt{nw\_wireless.critical\_curve()}.}\label{fig:nw:long-haul-wireless:critical}
\end{omfigure}

Availability follows from how often the critical rate is exceeded, which is a property of the climate along the route, published in rain-rate
statistics that this chapter does not use. With a labelled model in which it rains 5\% of the time with a mean rate of
\qty{4}{\milli\metre\per\hour} (exponentially distributed), one 59-kilometre hop at \qty{11}{\giga\hertz} soaked end to end is down 0.047\% of
the year, about four hours; twenty such hops failing independently would put the route on fibre for about 0.9\% of the time. The model is crude;
the structure is not: availability multiplies along a chain, and a route is only as dry as its wettest hop.

\section{Millimetre waves and lasers}

\begin{definition}[Millimetre-wave link, free-space optical link]\label{def:nw:long-haul-wireless:mmw}
A \emph{millimetre-wave link}\index{millimetre-wave link} is a radio link at frequencies of about 30 to 300~GHz (wavelengths of 10 to 1
millimetres), with wide channels and short hops. A \emph{free-space optical link}\index{free-space optical link} (FSO) carries data on a laser
beam through the air between two terminals; fog and heavy rain block it, and it needs a clear, stable line of sight.
\end{definition}

Inside a metropolitan area the trade-off reverses: hops are short, rain has little path to act on, and bandwidth matters. ICE's New Jersey
routes between Mahwah and Secaucus, run with Anova Financial Networks, combine millimetre-wave and free-space optics over 5-gigabit radio
channels; between Mahwah and Carteret they use several 1-gigabit channels. The two technologies fail differently (rain for millimetre waves,
fog for lasers), so a route that carries both loses both only when rain and fog come together.

\section{Shortwave radio across oceans}

\begin{definition}[Skywave propagation]\label{def:nw:long-haul-wireless:sky}
\emph{Skywave propagation}\index{skywave propagation} carries radio waves of a few to a few tens of megahertz beyond the horizon by reflection from
the ionosphere, one or more times, so that a single link can span an ocean.
\end{definition}

The geodesic from Secaucus NY4 to Slough LD4 is \qty{5551}{\kilo\metre}: \qty{18.52}{\milli\second} of light in vacuum, against
\qty{27.07}{\milli\second} in straight glass and \qty{29.48}{\milli\second} for the fastest published subsea route of chapter~13. A single reflection
from a layer 300~kilometres high lengthens a 5\,551-kilometre path by only about 0.6\% in a flat-earth approximation, so a skywave link can
beat any cable by about ten milliseconds before its equipment, with a bandwidth far below a microwave channel's and an availability that follows the ionosphere, hour by hour.

\begin{dated}{2026-09}{Shortwave for trading}\label{dat:nw:long-haul-wireless:hf}
Trade press reported in July 2023 that high-frequency traders had used shortwave links for several years to send trading data between U.S. and
foreign exchanges under experimental authorisations, and that a Shortwave Modernization Coalition wanted the \qtyrange{2}{25}{\mega\hertz} band
opened to fixed, long-distance, non-voice use. The FCC sought comment on the coalition's petition (RM-11953) by public notice on 30 June 2023.
\end{dated}

\section{Licences, towers, bandwidth and what is sent}

A microwave route is public in a way a fibre route is not: in the United States its paths are licensed, and Laughlin, Aguirre and Grundfest used
the FCC's licence records to document at least fifteen Chicago--New York networks, their towers and their lengths. The spectrum is narrow:
private networks there use mostly 30- and 40-megahertz channels in the 6- and 11-gigahertz bands, at about 140 to 190~megabits a second per
channel in their account. Venue-sold services are narrower still: 1 to 10~megabits of private bandwidth on ICE's Toronto service, 10~megabits on
SIX's Zurich--Frankfurt network.

What is sent is therefore chosen: the few fields of the few instruments that move another market (an index future's best price, a trade), in
compact formats, with everything else on fibre. The radio's own latency matters: Laughlin and his co-authors describe radios on the market in
2010 that added tens to thousands of microseconds per hop, mostly in \omterm{def:nw:long-haul-fibre:fec}{forward error correction} and interleaving, and estimate a few to twenty
microseconds per hop for the networks then being licensed. A route of twenty hops multiplies each microsecond by twenty.

\section{A storm across the route}

The model runs a storm across the Aurora--Carteret route: twenty hops of \qty{59}{\kilo\metre} on the great circle (\qty{1180.9}{\kilo\metre}),
a rain cell \qty{30}{\kilo\metre} across moving across the line at \qty{48}{\kilo\metre\per\hour}, with its rate falling linearly from a peak at
its centre to zero at its edge. While every hop holds its margin, the route runs at the published \qty{3.982}{\milli\second}; when one fails, it
falls back to fibre, which the model takes as a path 20\% longer than the geodesic, \qty{6.911}{\milli\second}. Everything but the geometry,
the coefficients and the radio route's published latency is an assumption.

\begin{omfigure}
\begin{tikzpicture}
\begin{axis}[width=11cm,height=5.8cm,xmin=0,xmax=37,ymin=3.5,ymax=7.5,xlabel={minutes since the cell reached the route},
  ylabel={route latency (ms, one way)},tick label style={font=\footnotesize},label style={font=\small},grid=major,
  grid style={gray!25},legend columns=3,legend style={font=\footnotesize,at={(0.5,-0.25)},anchor=north,draw=none},table/col sep=comma]
\addplot[omDef,thick,const plot] table[x=minute,y=p50]{figdata/networks/14-long-haul-wireless/storm.csv};
\addplot[omThm,thick,dashed,const plot] table[x=minute,y=p100]{figdata/networks/14-long-haul-wireless/storm.csv};
\addplot[omMeth,thick,dotted,const plot] table[x=minute,y=p150]{figdata/networks/14-long-haul-wireless/storm.csv};
\legend{peak 50 mm/h,peak 100 mm/h,peak 150 mm/h}
\end{axis}
\end{tikzpicture}

\omcaption{Simulation, not a measurement: the Aurora--Carteret route's one-way latency, minute by minute, while a rain cell \qty{30}{\kilo\metre}
across crosses one of its \qty{11}{\giga\hertz} hops at \qty{48}{\kilo\metre\per\hour}, for three peak rain rates: the radio route's published
\qty{3.982}{\milli\second}, or fibre at an assumed \qty{6.911}{\milli\second} while any hop is down. Data: \texttt{fig\_wireless.py},
\texttt{nw\_wireless.storm()}.}\label{fig:nw:long-haul-wireless:storm}
\end{omfigure}

With a peak of \qty{100}{\milli\metre\per\hour} the route spends 21 of the storm's 38 minutes on fibre, \qty{2.93}{\milli\second} slower; with
\qty{50}{\milli\metre\per\hour}, 5 minutes; with \qty{150}{\milli\metre\per\hour}, 26. The same storms at \qty{6}{\giga\hertz} cost 0, 0 and 7
minutes; at \qty{18}{\giga\hertz}, 22, 29 and 31. For the firms that race, those minutes are the afternoons of Shkilko and Sokolov's study; for
their counterparties, they are when the market is kinder.

\omcode{code/firm/radiolink/firm_radiolink.py}{130}{150}{The storm: the cell's chord across the route, the wet length and rain rate on each hop
under it, and the route's latency, radio or fibre, minute by minute.}\label{lst:nw:long-haul-wireless:storm}

\begin{method}[Evaluating a radio route]\label{met:nw:long-haul-wireless:eval}
\begin{enumerate}
\item From the licence records or the provider, get the towers: number of hops, lengths, frequencies. Compute the path's \omterm{def:nw:the-north-american-map:factor}{route factor}
  (chapter~10).
\item For each hop, compute the geometry and, with the provider's equipment figures, the \omterm{def:nw:long-haul-wireless:budget}{fade margin}; find the hops with the least critical
  rain rate.
\item Combine with the climate along the route to estimate availability; ask the provider for its measured availability and its outage history.
\item Price the fibre back-up, and design the switch-over: who detects the fade, how fast traffic moves, and what the strategies do meanwhile.
\item Decide what the radio carries: the few messages whose speed pays for it.
\end{enumerate}
\end{method}

\section{Tutorial: a hop, the rain and a storm}\label{tut:nw:long-haul-wireless:tut}

\begin{tutorial}
\textbf{Goal.} Compute a hop's geometry and budget, its rain limits by frequency, and a storm's cost in minutes on fibre.
\textbf{End state:} \cref{fig:nw:long-haul-wireless:gamma,fig:nw:long-haul-wireless:critical,fig:nw:long-haul-wireless:storm} and
\cref{tab:nw:long-haul-wireless:geometry,tab:nw:long-haul-wireless:budget}.
\begin{tutsteps}
\item \textbf{Geometry.} \texttt{firm\_radiolink.bulge\_m}, \texttt{fresnel\_m} and \texttt{tower\_m} (\cref{prop:nw:long-haul-wireless:geometry});
  \texttt{nw\_wireless.geometry\_rows()} fills \cref{tab:nw:long-haul-wireless:geometry}.
\item \textbf{Budget.} \texttt{Hop}, \texttt{rx\_dbm} and \texttt{fade\_margin\_db}; \texttt{budget\_rows()} fills
  \cref{tab:nw:long-haul-wireless:budget}.
\item \textbf{Rain.} The P.838-3 coefficients are data (\texttt{data/p838\_3.csv}), each row with its ledger reference;
  \texttt{critical\_rain} (\cref{lst:nw:long-haul-wireless:rain}) implements \cref{prop:nw:long-haul-wireless:critical}.
\item \textbf{The route and the storm.} \texttt{chain} places twenty hops on the great circle; \texttt{storm\_timeline}
  (\cref{lst:nw:long-haul-wireless:storm}) runs the cell across one of them.
\end{tutsteps}
\textbf{What to change next.} Use vertical polarisation (lower attenuation at these frequencies); shorten the hops to 40~kilometres and add
towers; move the storm along the route instead of across it.
\end{tutorial}

\section{Build: the radio link model}\label{bld:nw:long-haul-wireless:radiolink}

\begin{build}
\bfield{Purpose} Radio hops and chains with their geometry, budget and rain, and the failover to fibre: the air half of the routes that
chapters~15, 19 and~29 price and plan.
\bfield{Interface} \texttt{firm\_radiolink}: \texttt{bulge\_m}, \texttt{fresnel\_m}, \texttt{tower\_m}, \texttt{fspl\_db}, \texttt{dish\_gain\_dbi},
\texttt{Hop}, \texttt{rx\_dbm}, \texttt{fade\_margin\_db}, \texttt{coefficients}, \texttt{gamma\_db\_km}, \texttt{rain\_db}, \texttt{critical\_rain},
\texttt{chain}, \texttt{hop\_lengths}, \texttt{Storm}, \texttt{storm\_timeline}, \texttt{availability}.
\bfield{Rules} Rain coefficients come from the data file with their source, never typed in code; equipment figures are parameters; a storm is a
simulation and says so.
\bfield{Acceptance tests} \texttt{code/firm/radiolink/tests/}: the mid-hop rules, the free-space loss, a budget and the critical rain by hand, the
great-circle chain's end points, a dry and a wet storm, and availability at the extremes.
\bfield{Stretch} Terrain profiles along each hop; the ITU-R method for the effective path length of rain; correlated rain between neighbouring
hops.
\end{build}

\begin{omsources}
\item A.~Shkilko and K.~Sokolov, ``Every cloud has a silver lining: fast trading, microwave connectivity, and trading costs'', \emph{Journal of
Finance} 75(6) (2020).
\item Recommendation ITU-R P.838-3 (2005), specific attenuation model for rain.
\item G.~Laughlin, A.~Aguirre and J.~Grundfest (2014), section on microwave networks; ICE, New Jersey metro wireless routes; FCC Order DA 23-689
(2023) and Radio World (July 2023) on shortwave.
\end{omsources}

\section{Exercises}

\begin{exercise}[$\star$]\label{exo:nw:long-haul-wireless:1}
What are the \omterm{def:nw:long-haul-wireless:fresnel}{earth bulge} and the first Fresnel radius at the middle of a 40-kilometre hop at \qty{11}{\giga\hertz}, with $K = 4/3$?
\end{exercise}

\begin{exercise}[$\star$]\label{exo:nw:long-haul-wireless:2}
What is the specific attenuation of rain at \qty{25}{\milli\metre\per\hour} at 6 and at \qty{18}{\giga\hertz}?
\end{exercise}

\begin{exercise}[$\star$]\label{exo:nw:long-haul-wireless:3}
Why can a microwave network have a \omterm{def:nw:the-north-american-map:factor}{route factor} near 1.01 when a fibre route struggles to get below 1.2?
\end{exercise}

\begin{exercise}[$\star\star$]\label{exo:nw:long-haul-wireless:4}
Doubling a hop's length adds how many decibels of free-space loss, and what does it do to the critical rain rate if the rain covers the whole hop?
\end{exercise}

\begin{exercise}[$\star\star$]\label{exo:nw:long-haul-wireless:5}
Why do metropolitan routes pair millimetre waves with lasers?
\end{exercise}

\begin{exercise}[$\star\star$]\label{exo:nw:long-haul-wireless:6}
A shortwave link crosses the Atlantic ten milliseconds faster than any cable. What can it carry, and what can it not?
\end{exercise}

\begin{exercise}[$\star\star\star$]\label{exo:nw:long-haul-wireless:7}
\emph{Coding.} With \texttt{firm.radiolink}, find the peak rain rate at which the model's storm first puts the \qty{11}{\giga\hertz} route on fibre.
\end{exercise}

\begin{exercise}[$\star\star\star$]\label{exo:nw:long-haul-wireless:8}
\emph{Find the flaw.} ``Our route has twenty hops, each available 99.95\% of the year, so the route is available 99.95\% of the year.''
\end{exercise}

\section{Problem: The Storm Between Chicago and New Jersey}

\begin{problem}[{Weekend problem --- what a storm costs a radio route}]\label{pb:nw:long-haul-wireless:1}
A firm's radio route from Aurora to Carteret has twenty hops of about \qty{59}{\kilo\metre} at \qty{11}{\giga\hertz}, runs at \qty{3.982}{\milli\second}
one way, and falls back to fibre at \qty{6.911}{\milli\second}. Use the model's equipment and storm.

\medskip\noindent\textbf{Part I --- One hop.}
\begin{enumerate}
\item What are the mid-hop \omterm{def:nw:long-haul-wireless:fresnel}{earth bulge} and Fresnel radius, and the tower height for a 10-metre obstacle?
\item What are the hop's free-space loss, received power and \omterm{def:nw:long-haul-wireless:budget}{fade margin}?
\item At what rain rate does it lose its margin if the rain covers the whole hop?
\item And if the rain covers ten kilometres of it?
\end{enumerate}

\medskip\noindent\textbf{Part II --- The storm.}
\begin{enumerate}[resume]
\item How long does a cell 30 kilometres across take to cross the route at 48 kilometres an hour?
\item With a peak of \qty{100}{\milli\metre\per\hour}, how many minutes is the route on fibre?
\item How much slower is it during those minutes?
\item How does the answer change at 6 and at \qty{18}{\giga\hertz}?
\end{enumerate}

\medskip\noindent\textbf{Part III --- The year.}
\begin{enumerate}[resume]
\item With the labelled exceedance model, how long is one hop down in a year if the rain covers it end to end?
\item What does that give for twenty independent hops?
\item Why is the independence assumption doubtful?
\item What would shorter hops do?
\end{enumerate}

\medskip\noindent\textbf{Part IV --- The verdict.}
\begin{enumerate}[resume]
\item State the \emph{named result}: the rain rate at which a 59-kilometre hop at \qty{11}{\giga\hertz} loses its \omterm{def:nw:long-haul-wireless:budget}{fade margin}, and the minutes of
  fibre-only latency that a \qty{100}{\milli\metre\per\hour} storm crossing the route costs.
\item What did Shkilko and Sokolov find happens to trading costs during such episodes?
\item Why would a firm keep a 6-gigahertz route and an 18-gigahertz one?
\item What must the firm's systems do when the route switches to fibre?
\item Which of the model's inputs would you replace first with real data, and from where?
\item What does the storm change for a firm that does not race at all?
\item Is the fibre back-up's latency worth reducing? What would it take (chapter~13)?
\item In one sentence: what does the weather do to a latency race?
\end{enumerate}
\end{problem}

\section{Interview questions}

\begin{interviewq}[$\star$ \iqroles{developer}]\label{iq:nw:long-haul-wireless:1}
Why is microwave faster than fibre between two cities, and what does it cost in exchange?
\end{interviewq}

\begin{interviewq}[$\star\star$ \iqroles{developer}]\label{iq:nw:long-haul-wireless:2}
What limits the length of a microwave hop?
\end{interviewq}

\begin{interviewq}[$\star\star$ \iqroles{developer, trader}]\label{iq:nw:long-haul-wireless:3}
It is raining in Pennsylvania. What happens to your Chicago--New Jersey strategy, and what should your systems do?
\end{interviewq}

\begin{interviewq}[$\star\star$ \iqroles{trader, researcher}]\label{iq:nw:long-haul-wireless:4}
How would you use weather-driven outages of microwave networks to study the effect of speed on market quality?
\end{interviewq}

\begin{interviewq}[$\star\star$ \iqroles{developer}]\label{iq:nw:long-haul-wireless:5}
What would you send over a 10-megabit \omterm{def:nw:long-haul-wireless:mw}{microwave link} between two exchanges, and what would you leave on fibre?
\end{interviewq}

\begin{interviewq}[$\star\star\star$ \iqroles{developer, researcher}]\label{iq:nw:long-haul-wireless:6}
Compare microwave, millimetre-wave, lasers and shortwave for connecting two trading sites 50, 1\,000 and 6\,000 kilometres apart.
\end{interviewq}
